\scanpage{180}
\[
 f^\circ(\alpha_1,u_2(x_2),\ldots,u_n(x_n)).
\]
% slanted
\margin{\uline{NB} Si les \deleted{\unsure{Ext.}} $E_i$ sont munis de scindages, et leur structure de $\lambda$-triv. définie par des $b_i:F_i\mathbin{\unsure{-|\to}} E_i^\circ$, alors la $\lambda$-trivialisation de $*E_i$, scindé de \unsure{façon} can., est donnée par $\bigotimes b_i:\bigotimes F_i\mathbin{\unsure{-|\to}}\bigotimes E_i$.}

Ainsi, on a une opération interne $*$ dans la catégorie des extensions $\lambda$-scindées \inserted{(à quotients $F_i$ plats, au besoin)}, et en particulier dans celle des torseurs $\lambda$-scindés (car si les $F_i$ sont $\simeq k$, $\bigotimes F_i\simeq k$ canon.$^{\mathrm t}$).

\deleted{\unsure{Q}} Outre \unsure{sa|la} \unsure{donnée|condition} de commutativité, d'associativité (à vérifier leurs compatibilités !), a-t-on une unité ? Il faudrait pour le moins, \deleted{que} pour une unité, \unsure{q}
\diagram{\begin{tikzcd}[ampersand replacement=\&,column sep=small]
0\arrow[r,-]\& E^\circ\arrow[r,-]\& E\arrow[r,-]\& F\arrow[r]\& 0
\end{tikzcd}}
que $E^\circ$ \deleted{soit} $F$ soient unités pour le $\otimes$ ordinaire, donc \deleted{inversibles} \inserted{$\simeq k$}.
\diagram{\begin{tikzcd}[ampersand replacement=\&,column sep=small]
0\arrow[r]\&k\arrow[r]\&E\arrow[r]\&k\arrow[r]\&0
\end{tikzcd}}

\deleted{\unsure{Supp.|Une} tels} Choisissons scindage, \deleted{alors} la $\lambda$-triv. est définie par
\[
 \deleted{\unsure{k\to k}}\unsure{b_0|\beta_0|t}:k\longrightarrow k\qquad\text{i.e. }\beta_0\in k
\]

J'ai envie de prendre $\beta_0=1$. Prenons plus généralement $\beta$ quelconque, je dis qu'alors $E_\beta\mathbin{\underset{\lambda}{*}}E_2^{\unsure{+}}\simeq{}^{\beta}E_2$, muni de la $\lambda$ trivialisation $\beta u_2$ déduite de celle de $E_2$
% archive sheet number: 179
